% ============================================================
% PART I: mmWave Framework — summary slides 6-13
% ============================================================

\section{Part I: mmWave Framework}

% ----------------------------------------------------------
% SLIDE 6: Exact absorption law
% SAY: "Start from Poynting's theorem at a locally flat surface
%       point. The wave decomposes into TE and TM; Fresnel
%       gives the transmission for each. The ReLU gate is exact:
%       back-facing points receive zero power."
% DONT SAY: full Fresnel derivation (appendix)
% ----------------------------------------------------------
\begin{frame}{Exact absorption law}
\hypertarget{sec:part1}{}
\hypertarget{sum-p1-exact-law}{}

  \begin{columns}[T]
    \begin{column}{0.55\textwidth}
      Three factors determine local absorption:

      \begin{itemize}
        \item Illumination: incident power density $\Sinc$
        \item Material: effective Fresnel transmission $\Teff(\rr)$
        \item Geometry: visibility gate $\ReLU\!\bigl(\nhat\cdot(-\khat)\bigr)$
      \end{itemize}

      \vspace{6pt}
      \begin{resultbox}[Exact absorption law]
        \vspace{-6pt}
        \[
          \Sab(\rr) = \Sinc \;\Teff(\rr) \;\ReLU\!\bigl[\nhat(\rr)\cdot(-\khat)\bigr]
        \]
      \end{resultbox}

      \vspace{4pt}
      {\footnotesize
        $\Teff = |e_s|^2 T_s(\theta) + |e_p|^2 T_p(\theta)$: Fresnel-weighted average.\\
        $\ReLU$ gate is exact: back-facing points absorb nothing.\\
        \gotodetail{app-p1-fresnel}{Fresnel derivation}\;
        \gotodetail{app-p1-exact}{Exact law details}
      }
    \end{column}

    \begin{column}{0.4\textwidth}
      \begin{center}
        \includegraphics[width=\textwidth]{incidence_plane_fig.pdf}
      \end{center}
    \end{column}
  \end{columns}
\end{frame}

% ----------------------------------------------------------
% SLIDE 7: Pseudo-Brewster compensation
% SAY: "This is THE key physics insight. For biological tissue,
%       TE transmission decreases with angle, TM increases.
%       Their average stays near T_0. Azzam proved this in optics
%       for |n| > 2.5. Skin has |n| = 4.84. The material factor
%       collapses to a single number."
% ----------------------------------------------------------
\begin{frame}{Pseudo-Brewster compensation}
\hypertarget{sum-p1-brewster}{}

  \begin{columns}[T]
    \begin{column}{0.44\textwidth}
      $T_s$ decreases with angle; $T_p$ increases.\\
      Average $\Tavg(\theta)$ stays near $T_0$.

      \vspace{6pt}
      {\small
      \begin{tabular}{lccr}
        \toprule
        Tissue & $|\ntilde|$ & $T_0$ & $\Tavg/T_0$ var. \\
        \midrule
        Skin   & 4.84 & 0.54 & $\le 5.6\,\%$ \\
        Muscle & 5.62 & 0.48 & $\le 4.8\,\%$ \\
        Fat    & 2.05 & 0.77 & $\le 8.2\,\%$ \\
        Water  & 6.62 & 0.45 & $\le 3.9\,\%$ \\
        \bottomrule
      \end{tabular}
      }

      \vspace{6pt}
      \begin{resultbox}[Geometric absorption law]
        \vspace{-6pt}
        \[
          \Sab(\rr) \approx \Sinc \; T_0 \;\ReLU\!\bigl[\nhat(\rr)\cdot(-\khat)\bigr]
        \]
      \end{resultbox}
    \end{column}

    \begin{column}{0.52\textwidth}
      \begin{center}
        \includegraphics[width=\textwidth]{fresnel_curves.png}
      \end{center}
      {\footnotesize Azzam (2015): holds for $|\ntilde| > 2.5$.\\
        Biological tissue: $|\ntilde| \approx 3$--6. \checkmark\\
        \gotodetail{app-p1-pB}{Pseudo-Brewster details}}
    \end{column}
  \end{columns}
\end{frame}

% ----------------------------------------------------------
% SLIDE 8: Geometric consequences
% SAY: "Once the material collapses to T_0, everything is
%       geometry. Total absorbed power is T_0 times the shadow
%       area. Direction-average gives Cauchy's formula from 1841.
%       The self-shadowing correction eta is identical to ambient
%       occlusion in computer graphics."
% ----------------------------------------------------------
\begin{frame}{Geometric consequences}
\hypertarget{sum-p1-geometry}{}

  \begin{columns}[T]
    \begin{column}{0.52\textwidth}
      Total absorbed power: $P_{\mathrm{abs}} = \Sinc \, T_0 \, \Aperp(\khat)$\quad ($\Aperp$ = shadow area)

      \vspace{1pt}
      Self-shadowing on non-convex bodies:
      {\small $\displaystyle\eta(\rr) = \tfrac{1}{\pi}\!\int_{S^2}
        \ReLU[\nhat\!\cdot\!(-\khat)]\; O(\rr,\khat)\,\diff\Omega$}

      \vspace{1pt}
      $\equiv$ ambient occlusion (GPU graphics);
      Thelonious: $\Aab/A \approx 0.87$.

      \vspace{1pt}
      \begin{resultbox}[Generalised Cauchy formula]
        \vspace{-8pt}
        \[
          \langle P_{\mathrm{abs}} \rangle = \Sinc \, T_0 \, \frac{\Aab}{4}
        \]
        \vspace{-8pt}
      \end{resultbox}
      {\footnotesize Classical Cauchy ($\langle\Aperp\rangle = A/4$ for convex) is special case $\eta\equiv1$.\\
       \gotodetail{app-p1-geometry}{Cauchy proof \& ambient occlusion}}
    \end{column}

    \begin{column}{0.44\textwidth}
      \begin{center}
        \includegraphics[width=\textwidth,height=0.65\textheight,keepaspectratio]{eta_3d_phantom.png}
      \end{center}
      \vspace{-4pt}
      {\footnotesize Exposure fraction $\eta(\rr)$ on Thelonious.\\
       Red = exposed, blue = self-shadowed.}
    \end{column}
  \end{columns}
\end{frame}

% ----------------------------------------------------------
% SLIDE 9: Absorption directivity
% SAY: "The directivity D(k) encodes how the body's shape
%       modulates absorption. It ranges from 0.8 to 2.0 on
%       Thelonious. Smooth on the sphere: 49 SH coefficients
%       capture it to 1% RMS."
% ----------------------------------------------------------
\begin{frame}{Absorption directivity}

  \begin{columns}[T]
    \begin{column}{0.52\textwidth}
      \textbf{Factored result:}
      \[
        P_{\mathrm{abs}}(\khat) = \Sinc\, T_0\, \frac{\Aab}{4}\, D(\khat)
      \]

      \vspace{4pt}
      Separation of concerns:
      \begin{itemize}
        \item EM environment: $\Sinc$
        \item Tissue: $T_0$ (one number)
        \item Body geometry: $\Aab$ and $D(\khat)$, both precomputed
      \end{itemize}

      \vspace{6pt}
      $D(\khat)$ is smooth on $S^2$:\\
      SH expansion with $L \le 6$ (49 coeff.)\\
      gives RMS error $< 1\,\%$.
    \end{column}

    \begin{column}{0.44\textwidth}
      \begin{center}
        \includegraphics[width=\textwidth]{body_directivity_thelonious.png}
      \end{center}
      {\footnotesize Absorption directivity $D(\khat)$,\\
       Mollweide projection. Range: 0.8--2.0.}
    \end{column}
  \end{columns}
\end{frame}

% ----------------------------------------------------------
% SLIDE 10: Computational framework
% SAY: "For N paths from a ray tracer, the spatial dosimetry
%       map is a single matrix multiply with a ReLU. This is
%       literally a single-hidden-layer neural network.
%       The analogy is structural. The chain is differentiable."
% ----------------------------------------------------------
\begin{frame}{Computational framework}

  \textbf{Multi-source environment:} $N$ paths with directions $\KK$, power densities $\ssv$:

  \begin{resultbox}[Dosimetry as a ReLU neural network]
    \vspace{-6pt}
    \[
      \mathbf{S}_{\mathrm{ab}} = T_0 \;\ReLU(\NN\KK^\top)\; \ssv
    \]
  \end{resultbox}

  \vspace{4pt}
  \begin{columns}[T]
    \begin{column}{0.55\textwidth}
      \begin{itemize}
        \item $\NN$: surface normals (body mesh)
        \item $\KK$: path directions (ray tracer)
        \item $\ssv$: path power densities (ray tracer)
        \item ReLU: visibility gating (exact)
      \end{itemize}

      \vspace{4pt}
      \textbf{Differentiable}: gradients of SAR w.r.t.\ antenna positions propagate by backpropagation through any ray-tracing engine.
    \end{column}

    \begin{column}{0.4\textwidth}
      \textbf{Point source} ($d > 3\lambda$):
      \[
        P_{\mathrm{abs}} = P_t \, T_0 \, \frac{\Omega_{\mathrm{body}}}{4\pi}
      \]
      {\footnotesize Emissive view-factor formula\\
       of radiative heat transfer,\\
       with $T_0$ replacing emissivity.}
    \end{column}
  \end{columns}
\end{frame}

% ----------------------------------------------------------
% SLIDE 11: Validation
% SAY: "Mie theory is the exact wave-theoretic solution for
%       spheres. It tests both the Fresnel approximation and the
%       geometric-optics assumption. For body-scale objects at
%       mmWave, error is about 10%, dominated by diffraction into
%       the geometric shadow. The phantom test isolates the
%       Fresnel approximation: 0.35% total power error."
% DONT SAY: error decomposition details (appendix)
% ----------------------------------------------------------
\begin{frame}{Validation: Mie theory and phantom}
\hypertarget{sum-p1-validation}{}

  \begin{columns}[T]
    \begin{column}{0.48\textwidth}
      Mie theory (exact for lossy spheres):
      \begin{itemize}
        \item Tests Fresnel + geometric optics jointly
        \item Body-scale objects at mmWave: ${\sim}10\,\%$ error
        \item Dominated by diffraction into geometric shadow
        \item Error is conservative (underestimates)
      \end{itemize}

      \vspace{4pt}
      Thelonious phantom:
      \begin{itemize}
        \item Total absorbed power error: $0.35\,\%$
        \item Local RMS: $3.2\,\%$
        \item Isolates the Fresnel approximation from diffraction
      \end{itemize}
      {\footnotesize \gotodetail{app-p1-mie}{Mie details}\;
       \gotodetail{app-p1-3d}{Phantom results}}
    \end{column}

    \begin{column}{0.5\textwidth}
      \begin{center}
        \only<1>{%
          \includegraphics[width=\textwidth]{mie_panel_a.png}\\[2pt]
          {\footnotesize (a) Error vs size parameter at \SI{28}{\giga\hertz}.\\
           Body parts annotated by diameter.}
        }%
        \only<2>{%
          \includegraphics[width=\textwidth]{mie_panel_b.png}\\[2pt]
          {\footnotesize (b) Error vs frequency for different body parts.\\
           Torso at \SI{28}{\giga\hertz}: ${\sim}5\,\%$.}
        }%
        \only<3->{%
          \includegraphics[width=\textwidth]{mie_panel_c.png}\\[2pt]
          {\footnotesize (c) Fresnel asymptotic limit $R_{\mathrm{sphere}}(f)$.\\
           Conservative below \SI{40}{\giga\hertz}.}
        }%
      \end{center}
    \end{column}
  \end{columns}
\end{frame}

% ----------------------------------------------------------
% SLIDE 12: Error budget
% SAY: "The framework error is smaller than the tissue data
%       uncertainty. Fresnel approximation: 5% locally, 1%
%       integrated. Diffraction: 10% for body-part worst case.
%       Tissue dielectric data uncertainty: 20%. That is the
%       dominant source. Blue bars are conservative."
% ----------------------------------------------------------
\begin{frame}{Error budget}

  \begin{columns}[T]
    \begin{column}{0.42\textwidth}
      {\small
      \begin{tabular}{lr}
        \toprule
        Source & Magnitude \\
        \midrule
        Fresnel (local) & ${\sim}5\,\%$ \\
        Fresnel (integrated) & ${\sim}1\,\%$ \\
        Diffraction & ${\sim}10\,\%$\rlap{$^*$} \\
        Curvature & ${\sim}0.5\,\%$ \\
        \midrule
        Tissue uncertainty & ${\sim}20\,\%$ \\
        \midrule
        \textbf{Total framework} & \textbf{10--15\,\%} \\
        \bottomrule
      \end{tabular}
      }

      \vspace{2pt}
      {\footnotesize $^*$Worst case for body parts $< 10\lambda$.\\
       Conservative (underestimates absorption).\\[2pt]
       Framework error $<$ tissue data uncertainty.}
    \end{column}

    \begin{column}{0.54\textwidth}
      \begin{center}
        \only<1>{%
          \includegraphics[width=\textwidth]{error_budget_panel_a.png}\\[2pt]
          {\footnotesize (a) Best case: \SI{28}{\giga\hertz}, torso, multipath ($N\!=\!20$).\\
           All errors $\le 10\,\%$; dominated by diffraction and tissue.}
        }%
        \only<2->{%
          \includegraphics[width=\textwidth]{error_budget_panel_b.png}\\[2pt]
          {\footnotesize (b) Worst case: \SI{6}{\giga\hertz}, finger, single-path LOS.\\
           Near-field and diffraction dominate.}
        }%
      \end{center}
    \end{column}
  \end{columns}
\end{frame}

% ----------------------------------------------------------
% SLIDE 13: Part I detail index
% SAY: "These buttons jump to full derivations in the appendix.
%       We can use them during Q&A."
% ----------------------------------------------------------
\begin{frame}{Part I: detail index}
  \label{p1-index}

  \begin{columns}[T]
    \begin{column}{0.48\textwidth}
      \gotodetail{app-p1-setup}{Setup and notation}\\[4pt]
      \gotodetail{app-p1-flux}{Power flux derivation}\\[4pt]
      \gotodetail{app-p1-fresnel}{Fresnel coefficients}\\[4pt]
      \gotodetail{app-p1-exact}{Exact polarisation-aware law}\\[4pt]
      \gotodetail{app-p1-pB}{Pseudo-Brewster mechanism}\\[4pt]
      \gotodetail{app-p1-geometry}{Projected area and Cauchy}\\[4pt]
    \end{column}

    \begin{column}{0.48\textwidth}
      \gotodetail{app-p1-directivity}{Absorption directivity}\\[4pt]
      \gotodetail{app-p1-compliance}{Whole-body SAR compliance}\\[4pt]
      \gotodetail{app-p1-computation}{Computational framework}\\[4pt]
      \gotodetail{app-p1-corrections}{Higher-order corrections}\\[4pt]
      \gotodetail{app-p1-mie}{Mie validation details}\\[4pt]
      \gotodetail{app-p1-3d}{3D absorbed power density}\\[4pt]
      \gotodetail{app-p1-posture}{Posture sensitivity}\\[4pt]
      \gotodetail{app-p1-nearfield}{Near-field extension}\\[4pt]
      \gotodetail{app-p1-inter}{Inter-body reflections}\\[4pt]
    \end{column}
  \end{columns}
\end{frame}
